% ============================================================
% preamble.tex —— 共享导言区模板（通用版）
%
% 用途：由 md/PDF 转换而来的数学教材 LaTeX 项目的统一导言区，
%       提供页面设置、数学宏包、定理族环境、交叉引用与书名元信息。
%
% 说明：
%   * 文档类由 main.tex 声明（ctexbook），本文件只含宏包与配置，
%     因此不要在本文件中写 \documentclass。
%   * 编译方式（工作目录 = 项目根，即 main.tex 所在目录）：
%       xelatex -interaction=nonstopmode -output-directory=build main.tex
%     连续编译 2–3 遍，使目录、交叉引用（cleveref）与 PDF 书签收敛。
%   * 编码：所有文件均为 UTF-8。
%   * 本章节的文件命名约定：chapters/chNN.tex（如 ch01.tex）。
%   * 标签命名约定：前缀冒号格式，如 def:、thm:、lem:、prop:、
%     cor:、ex:、rem:（详见下方定理环境与 convert.py 的 id_map）。
% ============================================================

% ------------------------------------------------------------
% 〇、页面设置（收窄页边距，减少页面空白）
% ------------------------------------------------------------
% ctexbook 默认页边距偏大，版面留白过多。
% 以下收窄为：左右 2.2cm、上下 2.5cm，提升版面利用率。
% 如需再调，改动下列数值后重新编译两遍即可。
\usepackage{geometry}
\geometry{
  a4paper,
  left=2.2cm, right=2.2cm,   % 左右边距
  top=2.5cm, bottom=2.5cm,   % 上下边距
  headsep=0.6cm,             % 页眉与正文的距离
  footskip=0.9cm,            % 页码与正文底部的距离
}

% 允许 LaTeX 在紧急时拉伸字间距，避免长行内公式/长词溢出版心
% （配合收窄后的边距，消除 "Overfull \hbox" 警告）
\setlength{\emergencystretch}{3em}

% ------------------------------------------------------------
% 一、数学宏包
% ------------------------------------------------------------
% 注意加载顺序：先 amsmath / amssymb，再 amsthm（定理环境）与
% aliascnt（须在 amsthm 之后），再 mathtools（amsmath 增强），
% 最后 bm / mathrsfs / xcolor。
\usepackage{amsmath}     % 数学环境基础（equation、align 等）
\usepackage{amssymb}     % 数学符号（\varnothing、\lesssim 等）
\usepackage{amsthm}      % 定理环境与证明环境
\usepackage{aliascnt}    % 为共享编号序列的环境提供独立计数器名（供 cleveref 区分类型）
\usepackage{mathtools}   % amsmath 增强（\coloneqq、dcases 等）
\usepackage{bm}          % 粗体数学符号（\bm）
\usepackage{mathrsfs}    % 花写字母 \mathscr（OCR 识别的拓扑/代数符号）
\usepackage{xcolor}      % 文本颜色(\textcolor)—— 标记 OCR 损坏公式待校对

% 长公式允许适度跨页（教材排版常用）
\allowdisplaybreaks[1]

% ------------------------------------------------------------
% 二、图形
% ------------------------------------------------------------
\usepackage{graphicx}    % 插图
\graphicspath{{images/}} % 插图统一存放在项目 images/ 目录
\providecommand{\pandocbounded}[1]{#1}  % pandoc 3.x 图片包裹命令透传

% ------------------------------------------------------------
% 三、列表
% ------------------------------------------------------------
\usepackage{enumitem}    % 列表间距与标号格式微调
\setlist{itemsep=2pt, topsep=4pt, parsep=0pt}

% ------------------------------------------------------------
% 四、定理族环境（国内教材编号惯例）
% 计数器设计说明：
%   * theorem 为主计数器，按章编号（编号形如"定理 1.2"）；
%   * definition / lemma / proposition / corollary 需要与 theorem
%     "共享同一编号序列"，使"定义 1.1、定理 1.2、引理 1.3……"在同一
%     章内连续编号；但**不能直接共享同一个计数器**。
%     原因：cleveref 是依据 \label 绑定的"计数器名"来判定引用类型的。
%     若四者直接共享 theorem 计数器，它们的 label 全部绑定到 theorem
%     这一个计数器名上，于是无一例外被识别为定理类型，\cref 一律
%     渲染成"定理 X"；同时 \crefname{definition} 等条目永远不会被
%     查到，形同虚设。
%     方案：用 aliascnt 为四者各建立**独立计数器名**，同时保持与
%     theorem **共享同一编号序列**（编号仍然连续），从而既连续编号、
%     又能被 cleveref 正确区分类型。
%   * example、remark 使用独立计数器（同样按章编号），本身就彼此独立，
%     无需 aliascnt。
% 加载顺序约束：
%   * amsthm 须在 hyperref 之前加载；aliascnt 须在 amsthm 之后加载；
%   * 下列全部 \newtheorem 定义须在 cleveref 之前完成。
% ------------------------------------------------------------
% --- 4.1 声明主计数器（plain 风格：编号体斜体） ---
\newtheorem{theorem}{定理}[chapter]

% --- 4.2 共享编号序列、但各自拥有独立计数器名的环境 ---
% 每个环境三步：\newaliascnt 建独立计数器名（以 theorem 为别名）
%              -> \newtheorem 让环境使用该计数器
%              -> \aliascntresetthe 使新计数器与 theorem 同步计数
\theoremstyle{definition}                % 正文直体（黑体编号由 ctex 处理）
\newaliascnt{definition}{theorem}
\newtheorem{definition}[definition]{定义}
\aliascntresetthe{definition}

\theoremstyle{plain}                     % 正文斜体
\newaliascnt{lemma}{theorem}
\newtheorem{lemma}[lemma]{引理}
\aliascntresetthe{lemma}

\newaliascnt{proposition}{theorem}
\newtheorem{proposition}[proposition]{命题}
\aliascntresetthe{proposition}

\newaliascnt{corollary}{theorem}
\newtheorem{corollary}[corollary]{推论}
\aliascntresetthe{corollary}

% --- 4.3 独立计数器的环境（本就独立，无需 aliascnt） ---
\theoremstyle{definition}
\newtheorem{example}{例}[chapter]        % 例：definition 风格
\theoremstyle{remark}
\newtheorem{remark}{注}[chapter]         % 注：remark 风格

% --- 4.4 证明环境汉化 ---
\renewcommand{\proofname}{证}

% ------------------------------------------------------------
% 五、超链接与交叉引用
% ------------------------------------------------------------
\usepackage{hyperref}    % 必须在 cleveref 之前加载
\hypersetup{
  hidelinks,             % 链接不加彩框、不变色
  bookmarksnumbered,     % PDF 书签带章节编号
}

% cleveref 必须在 hyperref 之后加载
\usepackage{cleveref}
% 中文引用名（单数 / 复数取同一形式）
% 匹配规则：cleveref 通过 \label 所绑定的"计数器名"查表选取引用名，
% 而不是按环境名匹配。因此当多个环境共享同一计数器时，必须先用
% aliascnt 为其建立独立计数器名（见第 4.2 节），否则同组环境会共用
% 主计数器的引用名（全部显示"定理"），下面 definition / lemma /
% proposition / corollary 四条也会形同虚设、永不生效。
\crefname{theorem}{定理}{定理}
\crefname{definition}{定义}{定义}
\crefname{lemma}{引理}{引理}
\crefname{proposition}{命题}{命题}
\crefname{corollary}{推论}{推论}
\crefname{example}{例}{例}
\crefname{remark}{注}{注}
\crefname{figure}{图}{图}
\crefname{equation}{式}{式}
\crefname{chapter}{章}{章}

% ------------------------------------------------------------
% 六、书名元信息命令
% ------------------------------------------------------------
% main.tex 中调用：\bookinfo{书名}{作者}
% 同时写入 PDF 元数据（pdftitle / pdfauthor）。
% 提示：作者留空可避免标题页出现多余文字，即 \bookinfo{《书名》}{}。
% ------------------------------------------------------------
\newcommand{\bookinfo}[2]{%
  \title{#1}%
  \author{#2}%
  \date{}%
  \hypersetup{pdftitle={#1}, pdfauthor={#2}}%
}
